\begin{ouverture}[geometrie]{13}{Thalès et triangles semblables}{Partie 4 · Géométrie}
{Alignements · Parallèles · Réciproque · Échelles}
\objectif{Écrire les rapports de côtés correspondants.}
\objectif{Calculer et démontrer un parallélisme.}
\objectif{Reconnaître et utiliser des triangles semblables.}
\prerequis{Proportionnalité, fractions et produits en croix.}
\end{ouverture}
\pagestyle{otMain}\essentiel\fontsize{11}{15}\selectfont
\essSection{1}{Le théorème de Thalès}
\begin{propriete}Si A, M, B sont alignés, A, N, C sont alignés et $(MN)\parallel(BC)$, dans les configurations emboîtée ou en papillon, alors $\dfrac{AM}{AB}=\dfrac{AN}{AC}=\dfrac{MN}{BC}$. Les longueurs utilisées sont positives.\end{propriete}
\begin{savoirfaire}[Associer les côtés]Dans les triangles AMN et ABC, A correspond à A, M à B et N à C. Tous les numérateurs doivent appartenir au même triangle. Choisir ensuite le rapport contenant l’inconnue et des données connues.\end{savoirfaire}
\begin{exemple}Si $AM=3$, $AB=5$, $AC=7$ et les hypothèses sont vérifiées, $AN/7=3/5$, donc $AN=4{,}2$ dans la même unité.\end{exemple}
\begin{attention}Si M est entre A et B, $MB=AB-AM$. Dans une configuration en papillon, où A est entre M et B, $MB=AM+AB$. Lire l’ordre des points avant de calculer.\end{attention}
\clearpage
\essSection{2}{Réciproque et similitude}
\begin{propriete}[Réciproque de Thalès]Si A, M, B et A, N, C sont alignés dans le même ordre et si $AM/AB=AN/AC$, alors $(MN)\parallel(BC)$. Préciser les alignements et l’ordre avant de comparer les rapports.\end{propriete}
\begin{propriete}Si les alignements sont connus et si ces rapports diffèrent, les droites ne sont pas parallèles : sinon le théorème imposerait leur égalité.\end{propriete}
\begin{definition}Deux triangles sont semblables quand leurs angles sont égaux deux à deux. Deux angles égaux suffisent, car la somme des angles vaut $180^\circ$. Les côtés correspondants sont alors proportionnels.\end{definition}
\begin{propriete}Si le coefficient des longueurs est $k>0$, le rapport des périmètres est $k$ et celui des aires est $k^2$. Une correspondance erronée des sommets donne des rapports inutilisables.\end{propriete}
\begin{savoirfaire}[Rédiger]Énoncer les hypothèses, nommer le théorème, écrire l’égalité des rapports puis calculer. Dans une réciproque, conclure explicitement au parallélisme des deux droites nommées.\end{savoirfaire}
\clearpage\rubrique{Découverte}
\begin{blocLivre}{ex-13-1}\begin{exercice}[1]{Écrire les rapports}
Dans la figure, M appartient à [AB], N à [AC] et $(MN)\parallel(BC)$. Écrire les rapports égaux.
\label{fig-ex-13-1}\figureGeo{thales-01}
\end{exercice}\end{blocLivre}
\begin{blocLivre}{ex-13-2}\begin{exercice}[1]{Une longueur réduite}
M est sur [AB], N sur [AC] et $(MN)\parallel(BC)$. AM = 3 cm, AB = 5 cm, AC = 10 cm. Calculer AN.
\label{fig-ex-13-2}\figureGeo{thales-02}
\end{exercice}\end{blocLivre}
\begin{blocLivre}{ex-13-3}\begin{exercice}[1]{La petite base}
Dans la même configuration parallèle, AM = 4 cm, AB = 10 cm et BC = 15 cm. Calculer MN.
\label{fig-ex-13-3}\figureGeo{thales-03}
\end{exercice}\end{blocLivre}
\begin{blocLivre}{ex-13-4}\begin{exercice}[1]{Retrouver le grand côté}
M est sur [AB], N sur [AC], $(MN)\parallel(BC)$. AM = 6 cm, AN = 4 cm et AC = 10 cm. Calculer AB.
\label{fig-ex-13-4}\figureGeo{thales-04}
\end{exercice}\end{blocLivre}
\begin{blocLivre}{ex-13-5}\begin{exercice}[1]{De la partie au tout}
M est sur [AB], N sur [AC], $(MN)\parallel(BC)$. AM = 3 cm, MB = 2 cm et AC = 12 cm. Calculer AN.
\label{fig-ex-13-5}\figureGeo{thales-05}
\end{exercice}\end{blocLivre}

\clearpage\rubrique{Entraînement}
\begin{blocLivre}{ex-13-6}\begin{exercice}[2]{Démontrer un parallélisme}
A, M, B sont alignés dans cet ordre ; A, N, C aussi. AM = 4, AB = 10, AN = 6 et AC = 15, en cm. Montrer que MN et BC sont parallèles.
\label{fig-ex-13-6}\figureGeo{thales-06}
\end{exercice}\end{blocLivre}
\begin{blocLivre}{ex-13-7}\begin{exercice}[2]{Des droites non parallèles}
A, M, B et A, N, C sont alignés dans cet ordre. AM = 3, AB = 8, AN = 5, AC = 12, en cm. Les droites MN et BC sont-elles parallèles ?
\label{fig-ex-13-7}\figureGeo{thales-07}
\end{exercice}\end{blocLivre}
\begin{blocLivre}{ex-13-8}\begin{exercice}[2]{Le papillon}
M, A, B sont alignés dans cet ordre ; N, A, C aussi. $(MN)\parallel(BC)$, AM = 3 cm, AB = 5 cm et AC = 8 cm. Calculer AN.
\label{fig-ex-13-8}\figureGeo{thales-08}
\end{exercice}\end{blocLivre}
\begin{blocLivre}{ex-13-9}\begin{exercice}[2]{La distance entre deux points}
Dans le papillon précédent, AM = 3 cm et AB = 5 cm. Calculer MB. Si AN = 4,8 cm et AC = 8 cm, calculer NC.
\label{fig-ex-13-9}\figureGeo{thales-09}
\end{exercice}\end{blocLivre}
\begin{blocLivre}{ex-13-10}\begin{exercice}[2]{Des angles correspondants}
ABC et DEF ont respectivement les angles $(40^\circ,65^\circ,75^\circ)$ et $(65^\circ,75^\circ,40^\circ)$. Donner la correspondance des sommets. Si AB = 6 cm et FD = 9 cm, quel est le coefficient d’agrandissement ?
\label{fig-ex-13-10}\figureGeo{semblables-angles}
\end{exercice}\end{blocLivre}

\clearpage\rubrique{Problèmes}
\begin{blocLivre}{ex-13-11}\begin{exercice}[3]{Une aire agrandie}
Deux triangles semblables ont des côtés correspondants de 4 cm et 10 cm. L’aire du petit est 12 cm$^2$. Calculer le périmètre du grand si celui du petit est 15 cm, puis son aire.
\label{fig-ex-13-11}\figureGeo{semblables-aires}
\end{exercice}\end{blocLivre}
\begin{blocLivre}{ex-13-12}\begin{exercice}[3]{L’ombre d’un arbre}
Un bâton vertical de 1,5 m projette une ombre de 2 m. Au même instant, un arbre vertical projette une ombre de 12 m sur le même sol horizontal. Les rayons solaires sont parallèles. Calculer la hauteur de l’arbre.
\label{fig-ex-13-12}\figureGeo{thales-ombre}
\end{exercice}\end{blocLivre}
\begin{blocLivre}{ex-13-13}\begin{exercice}[3]{Une équation dans Thalès}
M est sur [AB], N sur [AC], MN est parallèle à BC. AM = 4 cm, MB = $x$ cm, AN = 6 cm et NC = 3 cm. Trouver $x$.
\label{fig-ex-13-13}\figureGeo{thales-13}
\end{exercice}\end{blocLivre}
\begin{blocLivre}{ex-13-14}\begin{exercice}[3]{Calculer une largeur inaccessible}
A, B, C forment un triangle. M est le milieu de AB et N le milieu de AC ; BC traverse une zone inaccessible. On mesure MN = 7,4 m. Justifier le parallélisme puis calculer BC.
\label{fig-ex-13-14}\figureGeo{thales-milieux}
\end{exercice}\end{blocLivre}
\begin{blocLivre}{ex-13-15}\begin{exercice}[3]{Un contrôle critique}
Dans une configuration emboîtée avec MN parallèle à BC, un élève écrit $AM/AB=NC/AC$. AM = 3 cm, AB = 8 cm, AC = 16 cm. Corriger et calculer NC.
\label{fig-ex-13-15}\figureGeo{thales-15}
\end{exercice}\end{blocLivre}

\clearpage
